\section{معادلات دو‌بعدی موج و گرما: مستطیل و دایره}
\def\EMfootchapter{فصل ۶ \ \lr{\textbullet}\ معادلات دو‌بعدی موج و گرما: مستطیل و دایره}
\def\EMtitlesize{26}\def\EMtitleskip{31}
\begin{frame}{فصل ۶: معادلات دو‌بعدی موج و گرما: مستطیل و دایره}
{\fontsize{20}{26}\selectfont\bfseries\color{SLdef}\EMalign مطالب این فصل\par}
\vspace{10bp}
\EMitem{\lr{\textbullet}}{\dimexpr\linewidth-26bp\relax}{معادلۀ موجی دو‌بعدی}
\EMitem{\lr{\textbullet}}{\dimexpr\linewidth-26bp\relax}{معادلۀ موجی دو‌بعدی با شرایط مرزی مستطیل}
\EMitem{\lr{\textbullet}}{\dimexpr\linewidth-26bp\relax}{معادلۀ گرمای دو‌بعدی با شرایط مرزی مستطیل}
\EMitem{\lr{\textbullet}}{\dimexpr\linewidth-26bp\relax}{بیان لاپلاسین در مختصات قطبی}
\EMitem{\lr{\textbullet}}{\dimexpr\linewidth-26bp\relax}{معادلۀ موجی دو‌بعدی با شرایط مرزی دایره}
\end{frame}

\def\EMtitlesize{26}\def\EMtitleskip{31}
\begin{frame}{معادلۀ موجی دو‌بعدی}
\begin{minipage}[t]{500bp}\vspace{0pt}
\begin{uncoverenv}<2->
\EMp{مدلسازی یک غشاء مرتعش (به‌دست آوردن معادلۀ حاکم بر \EMim{u(x,y,t)}).}
\end{uncoverenv}
\EMgap
\begin{uncoverenv}<3->
\EMp{\textbf{فرضیات فیزیکی:}}
\end{uncoverenv}
\EMgap
\begin{uncoverenv}<4->
\EMitem{\lr{\textbullet}}{\dimexpr\linewidth-26bp\relax}{جرم به صورت یکنواخت در طول غشاء توزیع شده (غشاء همگن)، کاملاً انعطاف‌پذیر است.}
\end{uncoverenv}
\EMgap
\begin{uncoverenv}<5->
\EMitem{\lr{\textbullet}}{\dimexpr\linewidth-26bp\relax}{غشاء چنان محکم کشیده شده‌است که نیروی جاذبه در آن قابل چشم‌پوشی است.}
\end{uncoverenv}
\EMgap
\begin{uncoverenv}<6->
\EMitem{\lr{\textbullet}}{\dimexpr\linewidth-26bp\relax}{میزان انحراف غشاء از حالت تعادل بسیار کوچک است.}
\end{uncoverenv}
\EMgap
\end{minipage}
\EMside{244.0bp}{\uncover<7->{\EMfigure{1.500}{m06-membrane}{المان \EMim{\Delta x\times\Delta y} از غشاء مرتعش و نیروهای کششی \EMim{T\Delta x} و \EMim{T\Delta y} وارد بر آن}}}
\end{frame}

\def\EMtitlesize{26}\def\EMtitleskip{31}
\begin{frame}{معادلۀ موجی دو‌بعدی (ادامه)}
\begin{uncoverenv}<2->
\EMp{با فرایند مدلسازی مشابه حالت یک‌بعدی، معادلۀ ارتعاش غشاء دوبعدی به دست می‌آید:}
\vspace{3bp}
\EMeq{\frac{\partial^2u}{\partial t^2} =c^2\left(\frac{\partial^2u}{\partial x^2}+\frac{\partial^2u}{\partial y^2}\right) =c^2\nabla^2u}
\end{uncoverenv}
\EMgap
\begin{uncoverenv}<3->
\EMp{برای حل یکتای این معادله نیاز به شرایط مرزی و شرایط اولیه داریم:}
\end{uncoverenv}
\EMgap
\begin{uncoverenv}<4->
\EMitem{\lr{\textbullet}}{\dimexpr\linewidth-26bp\relax}{شرایط مرزی: \EMim{u(x,y,t)=0} روی مرزهای غشاء}
\end{uncoverenv}
\EMgap
\begin{uncoverenv}<5->
\EMitem{\lr{\textbullet}}{\dimexpr\linewidth-26bp\relax}{شرایط اولیه: \EMim{u(x,y,0)=f(x,y)} و \EMim{u_t(x,y,0)=g(x,y)}، که در آن \EMim{u_t\triangleq\dfrac{\partial u}{\partial t}}}
\end{uncoverenv}
\end{frame}

\def\EMtitlesize{26}\def\EMtitleskip{31}
\begin{frame}{معادلۀ موجی دو‌بعدی با شرایط مرزی مستطیل}
\begin{minipage}[t]{500bp}\vspace{0pt}
\EMgap
\begin{uncoverenv}<3->
\EMeq{\frac{\partial^2u}{\partial t^2} =c^2\left(\frac{\partial^2u}{\partial x^2}+\frac{\partial^2u}{\partial y^2}\right)}
\end{uncoverenv}
\EMgap
\begin{uncoverenv}<4->
\EMitem{\lr{\textbullet}}{\dimexpr\linewidth-26bp\relax}{شرایط مرزی: \EMim{u(0,y,t)=u(a,y,t)=0} و \EMim{u(x,0,t)=u(x,b,t)=0}}
\end{uncoverenv}
\EMgap
\begin{uncoverenv}<5->
\EMitem{\lr{\textbullet}}{\dimexpr\linewidth-26bp\relax}{شرایط اولیه: \EMim{u(x,y,0)=f(x,y)} و \EMim{u_t(x,y,0)=g(x,y)}}
\end{uncoverenv}
\end{minipage}
\EMside{0.0bp}{\uncover<2->{\EMfigure{1.500}{m06-rectangle}{ناحیۀ مستطیلی \EMim{0\le x\le a}، \EMim{0\le y\le b}}}}
\end{frame}

\def\EMtitlesize{26}\def\EMtitleskip{31}
\begin{frame}{حل معادله}
\begin{uncoverenv}<2->
\EMp{در اینجا برای حل این معادله از روش \textbf{جداسازی متغیرها} استفاده می‌کنیم. فرض: \EMim{u(x,y,t)=X(x)\times Y(y)\times T(t)}.}
\vspace{3bp}
\EMeq{\frac{\partial^2u}{\partial x^2} =X''YT,\qquad \frac{\partial^2u}{\partial y^2} =XY''T,\qquad \frac{\partial^2u}{\partial t^2} =XYT''\quad \Longrightarrow\quad XYT'' =c^2\big(X''YT+XY''T\big)}
\end{uncoverenv}
\EMgap
\begin{uncoverenv}<3->
\EMeq{\frac{T''}{c^2T} =\frac{X''}X+\frac{Y''}Y =-\lambda^2\  \Rightarrow\ \frac{Y''}Y =-\lambda^2-\frac{X''}X =-\rho^2\  \Rightarrow\ \frac{X''}X =-\lambda^2+\rho^2 =-\mu^2}
\end{uncoverenv}
\EMgap
\begin{uncoverenv}<4->
\EMeq{\begin{cases}X''+\mu^2X=0\\ Y''+\rho^2Y=0\\ T''+c^2\lambda^2T=0\end{cases}\qquad \lambda^2 =\mu^2+\rho^2}
\end{uncoverenv}
\EMgap
\begin{uncoverenv}<5->
\EMp{\textbf{جواب در راستای \EMim{x}:}}
\vspace{3bp}
\EMeq{X(x) =A_1\cos\mu x+A_2\sin\mu x,\quad X(0) =0 \Rightarrow A_1 =0,\quad X(a) =0 \Rightarrow A_2\sin\mu a =0 \Rightarrow\mu a =n\pi \Rightarrow\mu_n \\[4pt] =\frac{n\pi}a}
\end{uncoverenv}
\end{frame}

\def\EMtitlesize{26}\def\EMtitleskip{31}
\begin{frame}{حل معادله (ادامه)}
\begin{uncoverenv}<2->
\EMeq{X_n(x) =A_n\sin\mu_nx,\qquad \mu_n =\frac{n\pi}a}
\end{uncoverenv}
\EMgap
\begin{uncoverenv}<3->
\EMp{\textbf{جواب در راستای \EMim{y}:}}
\vspace{3bp}
\EMeq{Y(y) =B_1\cos\rho y+B_2\sin\rho y,\quad Y(0) =0 \Rightarrow B_1 =0,\quad Y(b) =0 \Rightarrow B_2\sin\rho b =0 \Rightarrow\rho b =m\pi \Rightarrow\rho_m \\[4pt] =\frac{m\pi}b}
\end{uncoverenv}
\EMgap
\begin{uncoverenv}<4->
\EMeq{Y_m(y) =B_m\sin\rho_my,\qquad \rho_m =\frac{m\pi}b}
\end{uncoverenv}
\EMgap
\begin{uncoverenv}<5->
\EMp{\textbf{جواب در راستای زمان:}}
\vspace{3bp}
\EMeq{\lambda_{nm}^2 =\mu_n^2+\rho_m^2 =\left(\frac{n\pi}a\right)^2+\left(\frac{m\pi}b\right)^2,\qquad T_{nm}(t) =C_1\cos(c\lambda_{nm}t)+C_2\sin(c\lambda_{nm}t)}
\end{uncoverenv}
\EMgap
\begin{uncoverenv}<6->
\EMeq{u_{nm}(x,y,t) =\big[A_{nm}\cos(c\lambda_{nm}t)+B_{nm}\sin(c\lambda_{nm}t)\big]\sin\!\left(\frac{n\pi}ax\right)\sin\!\left(\frac{m\pi}by\right),\qquad n,m =1,2,\dots}
\end{uncoverenv}
\end{frame}

\def\EMtitlesize{26}\def\EMtitleskip{31}
\begin{frame}{حل معادله (ادامه)}
\begin{uncoverenv}<2->
\EMeq{u(x,y,t) =\sum_{n=1}^{\infty}\sum_{m=1}^{\infty}\big[A_{nm}\cos(c\lambda_{nm}t)+B_{nm}\sin(c\lambda_{nm}t)\big]\sin\!\left(\frac{n\pi}ax\right)\sin\!\left(\frac{m\pi}by\right)}
\end{uncoverenv}
\EMgap
\begin{uncoverenv}<3->
\EMp{\textbf{اعمال شرط اولیه \EMim{u(x,y,0)=f(x,y)}:}}
\vspace{3bp}
\EMeq{f(x,y) =\sum_{n=1}^{\infty}\sum_{m=1}^{\infty}A_{nm}\sin\!\left(\frac{n\pi}ax\right)\sin\!\left(\frac{m\pi}by\right) =\sum_{n=1}^{\infty}K_n(y)\sin\frac{n\pi}ax,\qquad K_n(y)\triangleq\sum_{m=1}^{\infty}A_{nm}\sin\frac{m\pi}by}
\end{uncoverenv}
\EMgap
\begin{uncoverenv}<4->
\EMeq{K_n(y) =\frac2a\int_0^af(x,y)\sin\frac{n\pi}ax\,dx,\qquad A_{nm} =\frac2b\int_0^bK_n(y)\sin\frac{m\pi}by\,dy}
\end{uncoverenv}
\EMgap
\begin{uncoverenv}<5->
\EMeq{A_{nm} =\frac4{ab}\int_0^b\left\{\int_0^af(x,y)\sin\frac{n\pi}ax\,dx\right\}\sin\frac{m\pi}by\,dy}
\end{uncoverenv}
\end{frame}

\def\EMtitlesize{26}\def\EMtitleskip{31}
\begin{frame}{حل معادله (ادامه)}
\begin{EMbox}{exe}{1}{}
\EMp{ثابت کنید:}
\vspace{3bp}
\EMeq{B_{nm} =\frac4{abc\lambda_{nm}}\int_0^b\left\{\int_0^ag(x,y)\sin\frac{n\pi}ax\,dx\right\}\sin\frac{m\pi}by\,dy}
\end{EMbox}
\begin{EMbox}{imp}{1}{}
\begin{uncoverenv}<2->
\EMeq{u(x,y,t) =\sum_{n=1}^{\infty}\sum_{m=1}^{\infty}\big[A_{nm}\cos(c\lambda_{nm}t)+B_{nm}\sin(c\lambda_{nm}t)\big]\sin\!\left(\frac{n\pi}ax\right)\sin\!\left(\frac{m\pi}by\right),\qquad \lambda_{nm}^2 \\[4pt] =\left(\frac{n\pi}a\right)^2+\left(\frac{m\pi}b\right)^2}
\end{uncoverenv}
\EMgap
\begin{uncoverenv}<3->
\EMeq{A_{nm} =\frac4{ab}\int_0^b\left\{\int_0^af(x,y)\sin\frac{n\pi}ax\,dx\right\}\sin\frac{m\pi}by\,dy,\qquad B_{nm} \\[4pt] =\frac4{abc\lambda_{nm}}\int_0^b\left\{\int_0^ag(x,y)\sin\frac{n\pi}ax\,dx\right\}\sin\frac{m\pi}by\,dy}
\end{uncoverenv}
\end{EMbox}
\end{frame}

\def\EMtitlesize{26}\def\EMtitleskip{31}
\begin{frame}{حل معادله (ادامه)}
\begin{EMbox}{ex}{1}{}
\EMp{مطلوب است جواب معادلۀ حرکت دو‌بعدی بالا با فرض \EMim{g(x,y)=0} و \EMim{f(x,y)=xy}.}
\EMgap
\begin{uncoverenv}<2->
\EMeq{g(x,y) =0\  \Rightarrow\ B_{nm} =0}
\end{uncoverenv}
\EMgap
\begin{uncoverenv}<3->
\EMeq{A_{nm} =\frac4{ab}\int_0^b\left\{\int_0^axy\sin\frac{n\pi}ax\,dx\right\}\sin\frac{m\pi}by\,dy =\frac4{ab}\left\{\left[\int_0^by\sin\frac{m\pi}by\,dy\right]\times\left[\int_0^ax\sin\frac{n\pi}ax\,dx\right]\right\}}
\end{uncoverenv}
\EMgap
\begin{uncoverenv}<4->
\EMeq{=\frac4{ab}\left[y\left(\frac{-b}{m\pi}\cos\frac{m\pi}by\right)-(1)\left(\frac{-b^2}{m^2\pi^2}\sin\frac{m\pi}by\right)\right]_0^b\times\left[x\left(\frac{-a}{n\pi}\cos\frac{n\pi}ax\right)-(1)\left(\frac{-a^2}{n^2\pi^2}\sin\frac{n\pi}ax\right)\right]_0^a}
\end{uncoverenv}
\EMgap
\begin{uncoverenv}<5->
\EMeq{=\frac4{ab}\cdot\frac{-b^2}{m\pi}(-1)^m\cdot\frac{-a^2}{n\pi}(-1)^n =\frac{4ab}{nm\pi^2}(-1)^{n+m}}
\end{uncoverenv}
\EMgap
\begin{uncoverenv}<6->
\EMeq{u(x,y,t) =\sum_{n=1}^{\infty}\sum_{m=1}^{\infty}\frac{4ab}{nm\pi^2}(-1)^{n+m}\cos(c\lambda_{nm}t)\sin\!\left(\frac{n\pi}ax\right)\sin\!\left(\frac{m\pi}by\right)}
\end{uncoverenv}
\end{EMbox}
\end{frame}

\def\EMtitlesize{26}\def\EMtitleskip{31}
\begin{frame}{معادلۀ گرمای دو‌بعدی با شرایط مرزی مستطیل}
\begin{uncoverenv}<2->
\EMeq{\frac{\partial u}{\partial t} =c^2\left(\frac{\partial^2u}{\partial x^2}+\frac{\partial^2u}{\partial y^2}\right) =c^2\nabla^2u}
\end{uncoverenv}
\EMgap
\begin{uncoverenv}<3->
\EMitem{\lr{\textbullet}}{\dimexpr\linewidth-26bp\relax}{شرط مرزی: \EMim{u(x,y,t)=0} روی مرزهای غشاء}
\end{uncoverenv}
\EMgap
\begin{uncoverenv}<4->
\EMitem{\lr{\textbullet}}{\dimexpr\linewidth-26bp\relax}{شرط اولیه: \EMim{u(x,y,0)=f(x,y)}}
\end{uncoverenv}
\end{frame}

\def\EMtitlesize{26}\def\EMtitleskip{31}
\begin{frame}{حل معادله}
\begin{uncoverenv}<2->
\EMp{فرض: \EMim{u(x,y,t)=X(x)\times Y(y)\times T(t)}.}
\vspace{3bp}
\EMeq{\frac{\partial^2u}{\partial x^2} =X''YT,\qquad \frac{\partial^2u}{\partial y^2} =XY''T,\qquad \frac{\partial u}{\partial t} =XYT'\quad \Longrightarrow\quad XYT' =c^2\big(X''YT+XY''T\big)}
\end{uncoverenv}
\EMgap
\begin{uncoverenv}<3->
\EMeq{\frac{T'}{c^2T} =\frac{X''}X+\frac{Y''}Y =-\lambda^2,\qquad \frac{Y''}Y =-\lambda^2-\frac{X''}X =-\rho^2,\qquad \frac{X''}X =-\lambda^2+\rho^2 =-\mu^2}
\end{uncoverenv}
\EMgap
\begin{uncoverenv}<4->
\EMeq{\begin{cases}X''+\mu^2X=0\\ Y''+\rho^2Y=0\\ T'+c^2\lambda^2T=0\end{cases}\qquad \lambda^2 =\mu^2+\rho^2}
\end{uncoverenv}
\EMgap
\begin{uncoverenv}<5->
\EMp{مشابه حالت موجی (\EMim{X_n(x)=A_n\sin\mu_nx} با \EMim{\mu_n=\frac{n\pi}a} و \EMim{Y_m(y)=B_m\sin\rho_my} با \EMim{\rho_m=\frac{m\pi}b}) و با}
\vspace{3bp}
\EMeq{T_{nm}(t) =C_1e^{-c^2\lambda_{nm}^2t},\qquad \lambda_{nm}^2 =\mu_n^2+\rho_m^2 =\left(\frac{n\pi}a\right)^2+\left(\frac{m\pi}b\right)^2}
\end{uncoverenv}
\end{frame}

\def\EMtitlesize{26}\def\EMtitleskip{31}
\begin{frame}{حل معادله (ادامه)}
\begin{uncoverenv}<2->
\EMeq{u_{nm}(x,y,t) =A_{nm}\sin\!\left(\frac{n\pi}ax\right)\sin\!\left(\frac{m\pi}by\right)e^{-c^2\lambda_{nm}^2t},\qquad n =1,2,\dots,\ m =1,2,\dots}
\end{uncoverenv}
\EMgap
\begin{uncoverenv}<3->
\EMeq{u(x,y,t) =\sum_{n=1}^{\infty}\sum_{m=1}^{\infty}A_{nm}\sin\!\left(\frac{n\pi}ax\right)\sin\!\left(\frac{m\pi}by\right)e^{-c^2\lambda_{nm}^2t}}
\end{uncoverenv}
\EMgap
\begin{uncoverenv}<4->
\EMp{\textbf{اعمال شرط اولیه:} دقیقاً مانند حالت موجی}
\vspace{3bp}
\EMeq{A_{nm} =\frac4{ab}\int_0^b\left\{\int_0^af(x,y)\sin\frac{n\pi}ax\,dx\right\}\sin\frac{m\pi}by\,dy}
\end{uncoverenv}
\end{frame}

\def\EMtitlesize{26}\def\EMtitleskip{31}
\begin{frame}{بیان لاپلاسین در مختصات قطبی}
\begin{minipage}[t]{500bp}\vspace{0pt}
\begin{uncoverenv}<2->
\EMp{هدف آن است که \EMim{\nabla^2u=u_{xx}+u_{yy}} در مختصات قطبی بر حسب \EMim{r,\theta} نوشته شود.}
\end{uncoverenv}
\EMgap
\EMgap
\begin{uncoverenv}<4->
\EMeq{r =\sqrt{x^2+y^2},\quad x =r\cos\theta,\qquad \theta =\tan^{-1}\!\left(\frac yx\right),\quad y \\[4pt] =r\sin\theta}
\end{uncoverenv}
\EMgap
\begin{uncoverenv}<5->
\EMeq{u_x =u_rr_x+u_\theta\theta_x\  \Rightarrow\ u_{xx} =\big(u_rr_x+u_\theta\theta_x\big)_x \\[4pt] =u_{rr}r_x^2+u_{\theta\theta}\theta_x^2+2r_x\theta_xu_{r\theta}+r_{xx}u_r+\theta_{xx}u_\theta}
\end{uncoverenv}
\EMgap
\begin{uncoverenv}<6->
\EMeq{u_y =u_rr_y+u_\theta\theta_y\  \Rightarrow\ u_{yy} \\[4pt] =u_{rr}r_y^2+u_{\theta\theta}\theta_y^2+2r_y\theta_yu_{r\theta}+r_{yy}u_r+\theta_{yy}u_\theta}
\end{uncoverenv}
\end{minipage}
\EMside{67.9bp}{\uncover<3->{\EMfigure{1.500}{m06-polar}{مختصات قطبی: \EMim{x=r\cos\theta}، \EMim{y=r\sin\theta}}}}
\end{frame}

\def\EMtitlesize{26}\def\EMtitleskip{31}
\begin{frame}{بیان لاپلاسین در مختصات قطبی (ادامه)}
\begin{uncoverenv}<2->
\EMeq{r_x =\frac xr,\quad r_{xx} =\frac{y^2}{r^3},\quad \theta_x =\frac{-y}{r^2},\quad \theta_{xx} =\frac{2xy}{r^4},\qquad r_y =\frac yr,\quad r_{yy} =\frac{x^2}{r^3},\quad \theta_y =\frac x{r^2},\quad \theta_{yy} =\frac{-2xy}{r^4}}
\end{uncoverenv}
\EMgap
\begin{uncoverenv}<3->
\EMeq{u_{xx} =\frac{x^2}{r^2}u_{rr}+\frac{y^2}{r^4}u_{\theta\theta}+\frac{-2xy}{r^3}u_{r\theta}+\frac{y^2}{r^3}u_r+\frac{2xy}{r^4}u_\theta}
\end{uncoverenv}
\EMgap
\begin{uncoverenv}<4->
\EMeq{u_{yy} =\frac{y^2}{r^2}u_{rr}+\frac{x^2}{r^4}u_{\theta\theta}+\frac{2xy}{r^3}u_{r\theta}+\frac{x^2}{r^3}u_r+\frac{-2xy}{r^4}u_\theta}
\end{uncoverenv}
\EMgap
\begin{uncoverenv}<5->
\EMp{با جمع این دو رابطه:}
\end{uncoverenv}
\EMgap
\begin{EMbox}{imp}{5}{}
\begin{uncoverenv}<6->
\EMeq{\nabla^2u =u_{rr}+\frac1{r^2}u_{\theta\theta}+\frac1ru_r}
\end{uncoverenv}
\end{EMbox}
\end{frame}

\def\EMtitlesize{26}\def\EMtitleskip{31}
\begin{frame}{معادلۀ موجی دو‌بعدی با شرایط مرزی دایره}
\begin{minipage}[t]{500bp}\vspace{0pt}
\EMgap
\begin{uncoverenv}<3->
\EMeq{\frac{\partial^2u}{\partial t^2} =c^2\nabla^2u,\qquad u(R,t) =0\ \ \EMt{(شرط مرزی)},\qquad u(r,0) \\[4pt] =f(r),\ \ u_t(r,0) =g(r)\ \ \EMt{(شرایط اولیه)}}
\end{uncoverenv}
\end{minipage}
\EMside{0.0bp}{\uncover<2->{\EMfigure{1.500}{m06-disc}{غشاء دایره‌ای به شعاع \EMim{R}}}}
\end{frame}

\def\EMtitlesize{26}\def\EMtitleskip{31}
\begin{frame}{حل معادله}
\begin{uncoverenv}<2->
\EMp{با توجه به تقارن دایروی در این مسئله، بهتر است آن را در مختصات قطبی حل کنیم. بر روی غشاء دایروی، ارتعاش تابعی از \EMim{\theta} نیست، لذا داریم:}
\vspace{3bp}
\EMeq{\begin{cases}\dfrac{\partial^2u}{\partial t^2}=c^2\left(u_{rr}+\dfrac1ru_r\right)\\\noalign{\vskip 2mm} u(R,t)=0,\quad u(r,0)=f(r),\quad u_t(r,0)=g(r)\end{cases}}
\end{uncoverenv}
\EMgap
\begin{uncoverenv}<3->
\EMp{حال با فرض تفکیک‌پذیری متغیرها \EMim{u(r,t)=P(r)\times T(t)} داریم:}
\vspace{3bp}
\EMeq{PT'' =c^2\left(P''T+\frac1rP'T\right)\  \Rightarrow\ \frac{T''}{c^2T} =\frac{P''}P+\frac1r\frac{P'}P =-\lambda^2\  \Rightarrow\ \begin{cases}P''+\dfrac1rP'+\lambda^2P=0\\\noalign{\vskip 1mm} T''+c^2\lambda^2T=0\end{cases}}
\end{uncoverenv}
\end{frame}

\def\EMtitlesize{26}\def\EMtitleskip{31}
\begin{frame}{معادلۀ دیفرانسیل بسل از مرتبۀ \EMim{n}}
\begin{minipage}[t]{500bp}\vspace{0pt}
\begin{uncoverenv}<2->
\EMeq{y''+\frac1xy'+\left(1-\frac{n^2}{x^2}\right)y =0\quad \Longrightarrow\quad y(x) \\[4pt] =AJ_n(x)+BY_n(x)}
\end{uncoverenv}
\EMgap
\begin{uncoverenv}<3->
\EMp{که در آن \EMim{J_n} توابع بسل نوع اول و از مرتبۀ \EMim{n} و \EMim{Y_n} توابع بسل نوع دوم و از مرتبۀ \EMim{n} هستند.}
\end{uncoverenv}
\EMgap
\end{minipage}
\EMside{174.0bp}{\uncover<4->{\EMfigure{0.880}{m06-bessel}{توابع بسل نوع اول \EMim{J_0,\dots,J_4} و نوع دوم \EMim{Y_0,\dots,Y_4} برای \EMim{0<x\le10}}}}
\end{frame}

\def\EMtitlesize{26}\def\EMtitleskip{31}
\begin{frame}{معادلۀ دیفرانسیل بسل از مرتبۀ \EMim{n} (ادامه)}
\begin{uncoverenv}<2->
\EMp{با تغییر متغیر \EMim{s=\lambda r}:}
\vspace{3bp}
\EMeq{P' =\frac{dP}{dr} =\frac{dP}{ds}\frac{ds}{dr} =\lambda\frac{dP}{ds},\qquad P'' =\lambda^2\frac{d^2P}{ds^2},\qquad \frac1r\lambda =\frac{\lambda^2}s}
\end{uncoverenv}
\EMgap
\begin{uncoverenv}<3->
\EMeq{\lambda^2\frac{d^2P}{ds^2}+\frac1r\lambda\frac{dP}{ds}+\lambda^2P =0\  \Rightarrow\ \frac{d^2P}{ds^2}+\frac1s\frac{dP}{ds}+P =0}
\end{uncoverenv}
\EMgap
\begin{uncoverenv}<4->
\EMp{این معادلۀ دیفرانسیل بسل از مرتبۀ صفر است (\EMim{n=0})، پس:}
\vspace{3bp}
\EMeq{P(r) =CJ_0(\lambda r)+DY_0(\lambda r)}
\end{uncoverenv}
\EMgap
\begin{uncoverenv}<5->
\EMp{\textbf{اعمال شرایط مرزی:}}
\end{uncoverenv}
\EMgap
\begin{uncoverenv}<6->
\EMitem{\lr{\textbullet}}{\dimexpr\linewidth-26bp\relax}{چون \EMim{P(0)\neq\infty} و \EMim{Y_0} در مبدأ بی‌کران است: \EMim{D=0} (ضریب \EMim{Y_0} صفر است) و \EMim{P(r)=CJ_0(\lambda r)}.}
\end{uncoverenv}
\end{frame}

\def\EMtitlesize{26}\def\EMtitleskip{31}
\begin{frame}{معادلۀ دیفرانسیل بسل از مرتبۀ \EMim{n} (ادامه)}
\begin{minipage}[t]{500bp}\vspace{0pt}
\begin{uncoverenv}<2->
\EMitem{\lr{\textbullet}}{\dimexpr\linewidth-26bp\relax}{\EMim{P(R)=0} نتیجه می‌دهد \EMim{J_0(\lambda R)=0}، یعنی \EMim{\lambda R=\alpha_m} که در آن \EMim{\alpha_m} ریشه‌های \EMim{J_0} هستند:}
\vspace{3bp}
\EMeq{\lambda =\lambda_m =\frac{\alpha_m}R,\quad m =1,2,\dots\quad \Longrightarrow\quad P_m(r) \\[4pt] =C_mJ_0\!\left(\frac{\alpha_m}Rr\right)}
\end{uncoverenv}
\EMgap
\EMgap
\begin{uncoverenv}<4->
\EMeq{T''+c^2\lambda_m^2T =0\  \Rightarrow\ T_m(t) =A_m\cos(c\lambda_mt)+B_m\sin(c\lambda_mt)}
\end{uncoverenv}
\EMgap
\begin{EMbox}{imp}{4}{}
\begin{uncoverenv}<5->
\EMeq{u(r,t) \\[4pt] =\sum_{m=1}^{\infty}\big(A_m\cos(c\lambda_mt)+B_m\sin(c\lambda_mt)\big)J_0\!\left(\frac{\alpha_m}Rr\right),\qquad \\[4pt] \lambda_m =\frac{\alpha_m}R}
\end{uncoverenv}
\end{EMbox}
\end{minipage}
\EMside{166.7bp}{\uncover<3->{\EMfigure{1.283}{m06-j0-zeros}{تابع \EMim{J_0(x)} و ریشه‌های آن \EMim{\alpha_1,\alpha_2,\alpha_3}}}}
\end{frame}

\def\EMtitlesize{26}\def\EMtitleskip{31}
\begin{frame}{معادلۀ دیفرانسیل بسل از مرتبۀ \EMim{n} (ادامه)}
\begin{EMbox}{rem}{1}{}
\begin{uncoverenv}<2->
\EMp{گزارشی کامل از معادلۀ دیفرانسیل بسل از مرتبۀ \EMim{n} و توابع بسل نوع اول و دوم و خواص آنها و نیز حل کامل معادلۀ دیفرانسیل موجی دوبعدی با شرایط مرزی دایره تهیه کنید (گزارش تایپ شده و در فرمت \lr{Word} باشد).}
\end{uncoverenv}
\end{EMbox}
\end{frame}
